\worksheettitle
  {M05\(\star\). Граница и внутренние отрезки}
  {\modulemark{M05\(\star\)} Усложнение}

\studentnameblock

\begin{notebox}[Главный вопрос]
  Дополнительные отрезки внутри фигуры могут делить её на части, но не
  добавляют новых сторон исходной внешней границе.
\end{notebox}

\quadrilateralregionsfigure

\begin{practicebox}[1. Исходная граница]
  Рассмотрите четырёхугольник $ABCD$ с внутренними отрезками.
  \begin{enumerate}[label=\alph*)]
    \item Сколько сторон осталось у внешней границы? \answerblank[15mm]
    \item Перечислите их. \answerblank[75mm]
    \item Почему точка $P$ не стала вершиной внешнего четырёхугольника?
      \writinglines{2}
  \end{enumerate}
\end{practicebox}

\begin{practicebox}[2. Считаем области]
  Сколько минимальных областей получилось внутри фигуры?
  \answerblank[18mm]

  Обведите границу каждой области на рисунке разным способом и запишите,
  сколько сторон имеет каждая область. \writinglines{3}
\end{practicebox}

\begin{checkpointbox}[Исследование]
  Нарисуйте пятиугольник и проведите внутри один отрезок между двумя его
  вершинами. Изменилось ли число сторон исходного пятиугольника? На сколько
  частей он мог разделиться? Объясните. \writinglines{4}
\end{checkpointbox}
